% Part: computability
% Chapter: computability-theory
% Section: introduction

\documentclass[../../../include/open-logic-section]{subfiles}

\begin{document}

\olfileid{cmp}{thy}{int}
\olsection{పరిచయం}

తర్కశాస్త్రంలో \emph{గణనీయతా సిద్ధాంతం}గా పిలిచే శాఖ ప్రమేయాలు,
సమితుల గణనీయత లేదా సాపేక్ష గణనీయతకు సంబంధించిన అంశాలను పరిశీలిస్తుంది.
క్లీనీ ప్రభావానికి నిదర్శనంగా, ఈ విషయాన్ని ఒకప్పుడు \emph{పునరావృత్త
సిద్ధాంతం} అని పిలిచేవారు; నేడు రెండు పేర్లూ సాధారణంగా వాడుకలో ఉన్నాయి.

ఏదో ఒక గణనా నమూనాలో గణించగల $f\colon \Nat \pto \Nat$ ప్రమేయాన్ని
\emph{పాక్షిక గణనీయ ప్రమేయం} అంటాం. $f$ సర్వనిర్వచితమైతే, $f$ను
కేవలం \emph{గణనీయ ప్రమేయం} అంటాం. గణనీయ లక్షణ ప్రమేయం~$\Char{R}$
గల సంబంధం~$R$ను కూడా గణనీయ సంబంధం అంటాం. $f$, $g$లు పాక్షిక
ప్రమేయాలైతే, $x$ వద్ద $f$ నిర్వచితమని, అంటే $x$ అనేది $f$ నిర్వచన
క్షేత్రంలో ఉందని సూచించడానికి $f(x) \fdefined$ అని రాస్తాం;
విరుద్ధంగా, $x$ వద్ద $f$ నిర్వచితం కాదని సూచించడానికి $f(x)
\fundefined$ అని రాస్తాం. $f(x)$, $g(x)$ రెండూ అనిర్వచితాలు, లేదా
రెండూ నిర్వచితమై సమానాలు అయినప్పుడు $f(x) \simeq g(x)$ అని రాస్తాం.

నిర్దిష్ట గణనా నమూనాను ప్రస్తావించకుండానే గణనీయతా సిద్ధాంతాన్ని
అధ్యయనం చేయవచ్చు. ఇందుకోసం సార్వత్రిక పాక్షిక గణనీయ
ప్రమేయం~$\fn{Un}(k, x)$ ఒకటి ఉందని చూపుతాం. దాని వల్ల పాక్షిక
గణనీయ ప్రమేయాలను లెక్కించవచ్చు. $k$వ ఏకస్థాన పాక్షిక గణనీయ ప్రమేయాన్ని
సూచించడానికి $\cfind{k}$ అనే సంకేతాన్ని తీసుకుంటాం; దాన్ని
$\cfind{k}(x) \simeq \fn{Un}(k, x)$గా నిర్వచిస్తాం. (క్లీనీ ఈ
ప్రయోజనానికి $\{ k \}$ వాడాడు; అయితే ఇటీవలి కాలంలో ఈ సంకేతం అంతగా
వాడబడటం లేదు.) ఇంకొంచెం సాధారణంగా, ఏ స్థానసంఖ్య గల పాక్షిక గణనీయ
ప్రమేయాలనైనా ఏకరీతిగా లెక్కించవచ్చు; $k$వ $n$-స్థాన పాక్షిక
పునరావృత్త ప్రమేయాన్ని సూచించడానికి $\cfind{k}[n]$ను వాడతాం.

$f(\vec x, y)$ సర్వనిర్వచిత లేదా పాక్షిక ప్రమేయమైతే,
$f(\vec x, y)=0$ అయ్యే అతి చిన్న~$y$ను తిరిగి ఇచ్చే $\vec x$ యొక్క
ప్రమేయమే $\umin{y}{f (\vec x, y)}$; ఇక్కడ $f(\vec x, 0)$, \dots,
$f(\vec x, y-1)$ అన్నీ నిర్వచితమని ఊహిస్తాం. అలాంటి $y$ లేకపోతే
$\umin{y}{f (\vec x, y)}$ అనిర్వచితం. $R(\vec x,y)$ ఒక సంబంధమైతే,
$R(\vec x,y)$ నిజమయ్యే అతి చిన్న~$y$గా $\umin{y}{R(\vec x,y)}$ను
నిర్వచిస్తాం; మరో మాటలో, $1 \tsub \Char{R}(\vec x,y)=0$ అయ్యే అతి
చిన్న~$y$ అది.

ఒక ప్రమేయం గణనీయమని చూపడానికి రెండు విధాలుగా ముందుకు వెళ్లవచ్చు:
\begin{enumerate}
\item కచ్చితంగా: ఒక ట్యూరింగ్ యంత్రాన్ని లేదా పాక్షిక పునరావృత్త
  ప్రమేయాన్ని స్పష్టంగా వర్ణించి, అది మనం ఉద్దేశించిన ప్రమేయాన్ని
  గణిస్తుందని చూపడం;
\item అనౌపచారికంగా: దాన్ని గణించే అల్గారిథంను వర్ణించి, చర్చ్
  సిద్ధాంతప్రతిపాదనను ఆశ్రయించడం.
\end{enumerate}
ఈ రెండింటి మధ్య ఖచ్చితమైన సరిహద్దు లేదు. ఒక అల్గారిథం యొక్క సమగ్ర
వివరణ, సూత్రప్రాయంగా సరైన ట్యూరింగ్ యంత్రాన్ని లేదా పాక్షిక పునరావృత్త
నిర్వచనాల అనుక్రమాన్ని ఎలా రూపొందించవచ్చో తగినంత స్పష్టం చేయాలి.
పూర్తిగా కచ్చితమైన నిర్వచనాలు సమాచారాత్మకంగా ఉండకపోవచ్చు; అందువల్ల ఈ
రెండు విధానాల మధ్య తగిన సమతుల్యతను వెతుకుతాం. సంక్షిప్తంగా చెప్పాలంటే,
కచ్చితమైన నిర్వచనాలను పొందవచ్చని చూపే, అంతర్బోధాత్మకమైనప్పటికీ కచ్చితమైన
నిరూపణలను ఇవ్వడానికి ప్రయత్నిస్తాం.

\end{document}
